\section{实验 Experiments}
\label{sec:experiments}

\subsection{设置}

我们在四种低资源语言上评估：阿姆哈拉语、孟加拉语、豪萨语与维吾尔语。
每种语言使用十万句单语语料。基线为直接沿用源语言词表，以及在目标语料上
从零训练的一元语言模型分词器~\cite{kudo2018subword}。

\begin{table}[t]
  \centering
  \small
  \caption{平均词元长度（字符数，越高越好）与下游命名实体识别 F1。}
  \label{tab:main}
  \begin{tabular}{lcccc}
    \toprule
    方法 & 阿姆哈拉 & 孟加拉 & 豪萨 & 维吾尔 \\
    \midrule
    沿用源词表   & 1.84 & 2.03 & 2.41 & 1.77 \\
    从零训练     & 2.29 & 2.51 & 2.88 & 2.16 \\
    本文方法     & \textbf{2.71} & \textbf{3.02} & \textbf{3.34} & \textbf{2.63} \\
    \midrule
    F1 提升      & +2.6 & +2.1 & +1.9 & +3.0 \\
    \bottomrule
  \end{tabular}
\end{table}

表~\ref{tab:main} 显示本文方法在四种语言上均优于两个基线。

\subsection{先验强度的影响}

\begin{figure}[t]
  \centering
  \begin{tikzpicture}
    \begin{axis}[
        width=0.86\linewidth, height=4.6cm,
        xlabel={先验强度 $\lambda$},
        ylabel={平均词元长度},
        xmin=0, xmax=1.0, ymin=2.0, ymax=3.2,
        grid=both, grid style={gray!20},
        tick label style={font=\scriptsize},
        label style={font=\scriptsize},
        legend style={font=\scriptsize,at={(0.98,0.04)},anchor=south east,draw=none,fill=none},
      ]
      \addplot[thick,mark=*,mark size=1.2pt] coordinates {
        (0,2.29) (0.1,2.48) (0.25,2.66) (0.5,2.71) (0.75,2.62) (1.0,2.41)
      };
      \addlegendentry{阿姆哈拉语}
      \addplot[thick,dashed,mark=square,mark size=1.2pt] coordinates {
        (0,2.16) (0.1,2.34) (0.25,2.57) (0.5,2.63) (0.75,2.55) (1.0,2.30)
      };
      \addlegendentry{维吾尔语}
    \end{axis}
  \end{tikzpicture}
  \caption{先验过强会压制目标语言自身的统计规律，$\lambda$ 超过 0.5 后性能下降。}
  \label{fig:lambda}
\end{figure}

图~\ref{fig:lambda} 表明先验强度存在一个明确的最优区间。

\subsection{失效情形}

该方法在两种情形下失效。第一，当源语言与目标语言的形态类型不同时，
例如源语言为屈折语而目标语言为黏着语，先验会引入错误约束，平均词元长度
反而下降百分之九。第二，当目标语料少于两万句时，式~\eqref{eq:objective}
的似然项过弱，结果几乎完全由先验决定，模型退化为源语言词表。
我们如实报告这两点，因为它们决定了该方法的适用边界。
